% TOP_PROBLEM: {"id": 30007008, "title": "Algebraicity of split Weil-type eightfold classes", "category_id": 5, "category": {"id": 5, "name": "algebraic_geometry", "display_name": "Algebraic Geometry", "description": "Geometric objects defined by polynomial equations.", "slug": "algebraic-geometry", "order_index": 5, "created_at": "2026-07-31T15:26:25.671Z"}, "status": "open"}
% FIRST_POSTED: 2026-10-06
% ATTEMPT_STATUS: solved
% OPENAI_CLAIM: {"schema_version":1,"source_repo":"https://github.com/openai/math","source_commit":"adc7f1241b42e322a6451854ab7e4b4c146bf78a","release_date":"2026-10-06","match":"full","resolution":"proved","independently_reviewed":false,"second_pass":"agreed","families":[{"family":"032","title":"Hodge and Kuga–Satake results for all projective K3 surfaces","subject":"Algebraic and complex geometry","manuscripts":[{"title":"Algebraicity of Weil classes on split abelian eightfolds","pdf_path":"preprints/Algebraicity-of-Weil-classes-on-split-abelian-eightfolds-September-18-2026/paper.pdf","readme_path":"preprints/Algebraicity-of-Weil-classes-on-split-abelian-eightfolds-September-18-2026/README.md","date":"2026-09-18","theorem_ref":"intro:main (sections/01-introduction.tex:51)","statement_sources":[{"path":"preprints/Algebraicity-of-Weil-classes-on-split-abelian-eightfolds-September-18-2026/build/sections/01-introduction.tex","line":51,"label":"intro:main"}]}],"lean":null,"reasoning_trace":null}]}
\documentclass[11pt]{article}
\usepackage[margin=25mm]{geometry}
\usepackage{fontspec}
\setmainfont{Latin Modern Roman}
\usepackage{amsmath,amssymb,mathtools}
\usepackage{unicode-math}
\setmathfont{Latin Modern Math}
\usepackage{enumitem}
\usepackage{xurl,hyperref,bookmark}
\hypersetup{colorlinks=true,urlcolor=blue}
\setcounter{secnumdepth}{0}
\setlength{\parindent}{0pt}
\setlength{\parskip}{5pt}
\setlength{\emergencystretch}{3em}
\begin{document}
\section*{Algebraicity of split Weil-type eightfold classes}

\subsection{OpenAI's claim}
Every rational Weil class of a split eight-dimensional abelian variety is claimed to have an algebraic representative, uniformly over imaginary-quadratic fields and compatible polarizations. The cited manuscript argues for that result; its relation to this problem is specified below.

\subsection{Mathematical statement of the claimed result}
\paragraph{Mathematical claim 1.} Let $K=\mathbb Q(\sqrt{-d})$ be imaginary quadratic, $A$ a complex abelian variety of dimension eight with a $K$-action, and $\lambda$ a compatible polarization of Weil type, with eigenspace multiplicities $(4,4)$. Put $V=H_1(A,\mathbb Q)$, let $\psi$ be the rational Riemann form and $J$ multiplication by $\sqrt{-d}$. Assume the associated Hermitian form
\[
h(x,y)=\psi(Jx,y)+\sqrt{-d}\,\psi(x,y)
\]
vanishes on a four-dimensional $K$-subspace of $V$ (the split condition). Define the rational Weil space $W_K(A)=\bigwedge_K^8H^1(A,\mathbb Q)\subset H^8(A,\mathbb Q)$ using the two conjugate $K$-eigenspace exterior powers. OpenAI claims
\[
W_K(A)\subset\operatorname{im}\bigl(\operatorname{CH}^4(A)_{\mathbb Q}\xrightarrow{\operatorname{cl}}H^8(A,\mathbb Q)\bigr).
\]
The field and compatible polarization type are arbitrary; extra endomorphisms are allowed. The statement concerns these Weil classes, rather than all Hodge classes on all varieties.

\subsection{Outline of the argument}
Manuscript 1 (family 032, 2026-09-18): An embedded spin Lagrangian detects the exceptional middle-cohomology component and makes two homological obstruction maps injective. Finite-cutoff weighted Floer operations are assembled into algebraic module identities; cancellation in alternating specialization then transports the desired Chern character to an algebraic cycle.

\subsection{Scope relative to this problem}
Manuscript 1 (family 032, 2026-09-18): The problem selects the mathematical statement summarized here, including its stated hypotheses and parameter range. The cited theorem gives a claimed proof of that assertion in full.

\subsection{Formal verification}
Family 032 has no Lean scope page in this snapshot; no formalization of these manuscript results is listed.

This site has not executed the Lean files or checked their dependencies. These descriptions report the scope of the linked declarations, rather than a new verification of the proofs.

\subsection{Status caveat}
OpenAI attributes these results to an internal model. Its release notice cautions that results lacking full formalization can contain mistakes. This site has not independently assessed the proofs; a solved label records the scope of the reported claim.

\subsection{Sources}
\begin{itemize}
\item OpenAI, \emph{Algebraicity of Weil classes on split abelian eightfolds}, 2026-09-18; citation key \texttt{OAI:Algebraicity-of-Weil-classes-on-split-abelian-eightfolds-September-18-2026}. \href{https://github.com/openai/math/blob/adc7f1241b42e322a6451854ab7e4b4c146bf78a/preprints/Algebraicity-of-Weil-classes-on-split-abelian-eightfolds-September-18-2026/paper.pdf}{PDF (GitHub)}; \href{https://github.com/openai/math/raw/adc7f1241b42e322a6451854ab7e4b4c146bf78a/preprints/Algebraicity-of-Weil-classes-on-split-abelian-eightfolds-September-18-2026/paper.pdf}{PDF (direct)}; \href{https://github.com/openai/math/blob/adc7f1241b42e322a6451854ab7e4b4c146bf78a/preprints/Algebraicity-of-Weil-classes-on-split-abelian-eightfolds-September-18-2026/README.md}{Manuscript page and citation}; \href{https://github.com/openai/math/blob/main/preprints/Algebraicity-of-Weil-classes-on-split-abelian-eightfolds-September-18-2026/paper.pdf}{Latest version}; Statement source: \href{https://github.com/openai/math/blob/adc7f1241b42e322a6451854ab7e4b4c146bf78a/preprints/Algebraicity-of-Weil-classes-on-split-abelian-eightfolds-September-18-2026/build/sections/01-introduction.tex\#L51}{intro:main}.
\item \href{https://github.com/openai/math/blob/adc7f1241b42e322a6451854ab7e4b4c146bf78a/CONTENTS.md}{OpenAI catalogue entry}.
\end{itemize}
\end{document}
